\section{Universal power ideals}
\label{sec:universal-power-v157}
The power-zero construction admits an equality of ideals before any
restriction to a pencil or to a valuation ring. This is the point at
which the reciprocal section acquires a global boundary geometry.
Let $V$ have dimension $n\geq2$, put $M=\Sym^2V$, and use the algebra
$B_{n,h}$ of Theorem~\ref{thm:quadratic-section-v155} with $W=V^*$.
Thus $(B_{n,h})_1=M$. If $A$ is a symmetric matrix with entries in a
commutative $\C$-algebra $T$, denote its element of $M\otimes T$ by
$\ell_A$, and set
\begin{equation}\label{eq:universal-power-definition-v157}
 J_p(A)=\im\bigl((B_{n,h})_p^*\otimes T\longrightarrow T,
                  \quad\varphi\longmapsto\varphi(\ell_A^p)\bigr).
\end{equation}
Write $I_q(A)$ for the ideal of all $q$-minors, with $I_0(A)=T$.
In this section powers of ideals are ordinary powers unless parentheses
explicitly indicate symbolic powers.

\begin{theorem}[Universal determinantal formula]
\label{thm:universal-power-v157}
For every commutative $\C$-algebra $T$, every symmetric $n\times n$
matrix $A$ over $T$, and $0\leq p\leq nh$, write $p=hq+s$ with
$0\leq s<h$. Then
\begin{equation}\label{eq:universal-power-v157}
             J_p(A)=I_q(A)^{h-s}I_{q+1}(A)^s.
\end{equation}
When $s=0$ the second factor is omitted. For $p>nh$, $J_p(A)=0$.
These are equalities of ideals, not of radicals or integral closures,
and commute with every homomorphism of $\C$-algebras. In particular,
for every symmetric map $A:T^n\to T^n$,
\begin{equation}\label{eq:universal-Fitting-v157}
       J_{h(n-j)}(A)=\Fitt_j(\coker A)^h,\qquad 0\leq j\leq n.
\end{equation}
The formulas also hold locally, and hence as ideal sheaves, for
symmetric maps of vector bundles with a line-bundle twist on an
arbitrary $\C$-scheme.
\end{theorem}

We separate the classical determinantal input. For an arbitrary matrix
$X$ over a characteristic-zero polynomial ring, the balancing inclusion
\begin{equation}\label{eq:minor-balancing-v157}
              I_u(X)I_v(X)\subseteq I_{u+1}(X)I_{v-1}(X),
              \qquad 0\leq u\leq v-2,
\end{equation}
is the determinantal symmetrization lemma. See
\cite[Lemma~2.5 and Theorem~2.4]{BrunsConca2003}, which attributes the
characteristic-zero product theorem to De Concini--Eisenbud--Procesi.
Because this is a polynomial ideal inclusion, it remains true after
specialization to a symmetric matrix; no flatness of that specialization
is needed. This classical lemma is not a claim of the present paper.

\begin{lemma}[The bounded-row generating space]
\label{lem:bounded-row-v157}
Let $S=\Sym(\Sym^2V^*)$ and let $A$ be its universal symmetric matrix.
For $p=hq+s$, the degree-$p$ part generating
$I_q(A)^{h-s}I_{q+1}(A)^s$ is
\begin{equation}\label{eq:bounded-row-v157}
 \bigoplus_{\substack{|\lambda|=p,\ \ell(\lambda)\leq n\\
                         \lambda_1\leq h}}\mathbb S_{2\lambda}(V^*)
                         \ \subseteq S_p.
\end{equation}
\end{lemma}
\begin{proof}
The symmetric Cauchy decomposition of $S_p$ is multiplicity free,
with summands $\mathbb S_{2\lambda}(V^*)$, $|\lambda|=p$; see
\cite[\S2.4]{HenriquesVarbaro2016}. We may use the dual general
linear group to regard these as polynomial representations. Every
minor has torus weight at most two in any one coordinate: its row
set and column set each use that coordinate at most once. A product
of $h$ minors therefore has no weight with a coordinate exceeding
$2h$. Its invariant span cannot contain a highest weight
$2\lambda$ with $\lambda_1>h$. This proves one inclusion.

For the converse, let $c_1\geq\cdots\geq c_h\geq0$ be the column
lengths of a partition $\lambda$ occurring on the right, padded with
zeros. The product of the leading principal minors of orders $c_i$
is a nonzero highest-weight vector of weight $2\lambda$ in $S_p$.
It belongs initially to $\prod_{i=1}^h I_{c_i}(A)$.
Whenever two column lengths differ by at least two, apply
\eqref{eq:minor-balancing-v157} to increase the smaller length by one
and decrease the larger by one. The sum of the squares of the lengths
strictly decreases, so the procedure terminates. At termination the
lengths are $q+1$ repeated $s$ times and $q$ repeated $h-s$ times,
since their sum is $p$. Thus that highest-weight vector belongs to
$I_q^{h-s}I_{q+1}^s$. The latter ideal is invariant; its degree-$p$
part contains the full irreducible summand. Multiplicity freeness
proves the equality. The case $q=0$ includes zero-length minors,
which are the constant $1$.
\end{proof}

\begin{proof}[Proof of Theorem~\ref{thm:universal-power-v157}]
First work over $S$ with its universal matrix. The coefficient map of
$\ell_A^p$ is the dual of the multiplication quotient
$\Sym^p M\twoheadrightarrow (B_{n,h})_p$, up to nonzero multinomial
constants. Consequently its degree-$p$ coefficient span is exactly
$(B_{n,h})_p^*\subseteq S_p$. The character formula
\eqref{eq:section-character-v155} identifies that span with
\eqref{eq:bounded-row-v157}. Both ideals in
\eqref{eq:universal-power-v157} are generated in degree $p$, so the
lemma proves their equality in $S$. This step uses the full coefficient
space, not just evaluation on diagonal matrices or closed points.

For a matrix over $T$, substitute its entries in this polynomial
identity. Taking the image of the fixed finite-dimensional coefficient
map commutes with scalar extension, as does taking minors and products
of ideals. This proves the assertion even when $T$ is nonreduced,
non-Noetherian, or not a domain. The Fitting formula follows from its
minor presentation. Changing a vector-bundle frame applies an invertible
congruence to $A$ and an invertible linear transformation to each
coefficient space; changing a local frame of the twisting line multiplies
generators by a unit. Therefore the local ideal identities glue. Finally,
$(B_{n,h})_p=0$ for $p>nh$.
\end{proof}

\begin{example}[A nonreduced test base]
\label{ex:nilpotent-base-v157}
Take $T=\C[\epsilon]/(\epsilon^3)$, $n=h=2$, and
$A=\diag(\epsilon,0)$. Then $J_1(A)=(\epsilon)$ and
$J_2(A)=(\epsilon^2)\ne0$, whereas $J_3(A)=J_4(A)=0$.
Every closed-point matrix is zero. Thus equality on geometric fibres
would not establish the theorem. Nor does
\eqref{eq:universal-Fitting-v157} license taking unique ideal roots:
for example, in $\C[\epsilon]/(\epsilon^2)$ the distinct ideals
$(\epsilon)$ and $0$ have the same square. At $h=1$ the power ideals
are the Fitting ideals themselves over every base.
\end{example}

\section{The reciprocal power graph and complete quadrics}
\label{sec:power-graph-v157}
Let $P=\Pj(M)$, let $P^{\mathrm{nd}}$ be the nonsingular-quadric open,
and define $\mathcal G_h$ as the scheme-theoretic closure of the graph
\begin{equation}\label{eq:power-graph-v157}
 P^{\mathrm{nd}}\longrightarrow
 P\times\prod_{p=1}^{nh-1}\Pj((B_{n,h})_p),
 \qquad [A]\longmapsto([A],[\ell_A],\ldots,[\ell_A^{nh-1}]).
\end{equation}
The first factor records the base point. We use the reduced, integral
graph closure, equivalently its scheme-theoretic image; no additional
normalization is inserted in this definition. Let $\mathrm{CQ}(V)$
be the classical space of complete quadrics: the graph closure of
$[A]\mapsto([\bigwedge^2A],\ldots,[\bigwedge^{n-1}A])$ over $P$.
For $n=2$ the empty product gives $\mathrm{CQ}(V)=P$.

\begin{theorem}[A compactification determined by multiplication]
\label{thm:power-graph-v157}
For every $h\geq1$ there is a canonical $\operatorname{GL}(V)$-equivariant
isomorphism over $P$
\begin{equation}\label{eq:power-graph-CQ-v157}
                      \mathcal G_h\simeq\mathrm{CQ}(V).
\end{equation}
In particular the graph is smooth and projective, is independent of
$h$ as a scheme over $P$, and has a simple-normal-crossing boundary.
Write $\pi:\mathrm{CQ}(V)\to P$, and let $E_r$ be its boundary divisor
of rank break $r$, $1\leq r<n$; for $r<n-1$ it is exceptional over
rank at most $r$, and $E_{n-1}$ is the strict transform of the determinant
divisor. Then every power ideal has the exact total transform
\begin{equation}\label{eq:power-boundary-v157}
 J_p\OO_{\mathrm{CQ}(V)}=
 \OO_{\mathrm{CQ}(V)}\left(-\sum_{r=1}^{n-1}(p-hr)_+E_r\right),
                    \qquad 0\leq p\leq nh,
\end{equation}
where $(a)_+=\max(a,0)$. Removing this common divisor from the universal
$p$th power gives a nowhere-zero projective section on the whole
compactification.
\end{theorem}

The smoothness, boundary and flag interpretation of complete quadrics
are classical, due to the complete-form construction; see
\cite{Vainsencher1982,Massarenti2020}. The assertion here identifies
that particular compactification, with all of its scheme structure,
from the multiplication in one reciprocal normalization-fibre section.
It does not assert that complete quadrics is a newly constructed variety.

\begin{proof}
We first recall the graph calculation with its scheme structure.
If rational maps from an integral scheme are given by generators of
nonzero ideals $K_1,\ldots,K_a$, their simultaneous graph closure is
$\operatorname{Bl}_{K_1\cdots K_a}$ of the base. Indeed, after the Segre
embedding the coordinates are all products of one generator from each
ideal. The closure of this graph is the Proj of the Rees algebra of
that product: on a generator chart its coordinate ring is the subring
of the function field obtained by adjoining the generator ratios.
These are exactly the affine charts of both constructions. Tensoring
an ideal by an invertible ideal, or replacing it by a positive power,
does not change the blow-up, the latter by the Veronese description
of Proj.

Apply this calculation on $P$. The sheaf $I_1\OO_P$ is the unit ideal,
and $I_n\OO_P$ is invertible. For each $2\leq q\leq n-1$, the exponent
of $I_q$ in $\prod_{p=1}^{nh-1}J_p$ is
\[
 \sum_{s=0}^{h-1}(h-s)+\sum_{s=1}^{h-1}s=h^2.
\]
Theorem~\ref{thm:universal-power-v157} therefore gives, up to an
invertible ideal,
\[
       \prod_{p=1}^{nh-1}J_p
                =\left(\prod_{q=2}^{n-1}I_q\right)^{h^2}.
\]
The same graph calculation for the exterior-power maps identifies
$\mathrm{CQ}(V)$ with the blow-up of $\prod_{q=2}^{n-1}I_q$.
This proves \eqref{eq:power-graph-CQ-v157} as an isomorphism of schemes,
not only after normalization. Both isomorphisms are the identity on
the common dense open, hence are canonical and equivariant.

For completeness, the local calculation behind the boundary formula
is as follows. The complete-quadric construction blows up the strict
rank loci successively. On its Gaussian charts, after a congruence
and projective normalization, its matrix has the form
\begin{equation}\label{eq:CQ-chart-v157}
 U\diag(1,\tau_1,\tau_1\tau_2,\ldots,
                         \tau_1\cdots\tau_{n-1})U^{\mathsf t},
\end{equation}
with $U$ lower unitriangular and $E_r=(\tau_r=0)$.
One obtains these charts inductively by making the first pivot one,
taking a Schur complement, blowing up the zero residual matrix, and
repeating. At each step the next nonzero residual quadratic form has
a nonisotropic vector, so congruence translates of these charts cover
the construction. This also explains the smooth charts and the normal
crossings stated in the classical construction.

Cauchy--Binet, applied to \eqref{eq:CQ-chart-v157}, shows that every
$q$-minor is divisible by
$\tau_1^{q-1}\tau_2^{q-2}\cdots\tau_{q-1}$.
The leading principal $q$-minor is exactly that product. Thus, on
these charts and therefore globally,
\begin{equation}\label{eq:minors-on-CQ-v157}
 I_q\OO_{\mathrm{CQ}(V)}=
              \OO\left(-\sum_{r<q}(q-r)E_r\right).
\end{equation}
Using \eqref{eq:universal-power-v157}, its exponent along $E_r$ is
$(h-s)(q-r)_++s(q+1-r)_+=(p-hr)_+$.
This proves \eqref{eq:power-boundary-v157}. The coefficient sections
generate the transformed ideal; division by its local generator makes
them generate the unit ideal, proving the last assertion.
\end{proof}

\begin{corollary}[Native boundary strata]
\label{cor:native-strata-v157}
The power graph has one congruence orbit for every subset
$I\subseteq\{1,\ldots,n-1\}$. It is the stratum
$E_I^\circ=(\bigcap_{r\in I}E_r)\setminus\bigcup_{r\notin I}E_r$,
of codimension $|I|$, and its closure is $\bigcap_{r\in I}E_r$.
Its points are flags with nondegenerate projective quadratic forms on
successive quotients; the quotient dimensions are the successive gaps
in $\{0\}\cup I\cup\{n\}$. The simultaneous normalized power values
recover this complete-quadric datum. They are not the one common limit
for the full cotangent-coordinate group.
\end{corollary}
\begin{proof}
The flag description and boundary stratification are the classical
complete-quadric description. The general linear group acts transitively
on flags of fixed dimensions, and over $\C$ every nondegenerate
quadratic form on a fixed-dimensional quotient is congruent to the
identity. Independent diagonal scalings remove the projective scalars.
The chart \eqref{eq:CQ-chart-v157} realizes each subset of vanishing
$\tau_r$ and permits further coordinates to vanish, giving the stated
codimensions and closures. The graph is a closed subscheme of the
product in \eqref{eq:power-graph-v157}; its projections therefore
retain the point, including its flag and quotient forms. The group
here is the native congruence group of $V$.
\end{proof}

\begin{remark}[Relative construction and base change]
\label{rem:relative-CQ-v157}
For a vector bundle $\mathcal V$ on any $\C$-scheme $S$, the associated
bundle with fibre $\mathrm{CQ}(\C^n)$ is a relative smooth projective
scheme over $S$. Equivariance glues the power maps and the boundary
formula. Formation of this associated bundle and the polynomial ideal
identity commute with arbitrary base change. This does not say that
taking the graph closure of an arbitrary restricted family commutes
with nonflat base change. The relative ambient resolution and the
closure of an individual family's graph are different constructions.
\end{remark}

\section{Boundary contacts and spectral schemes}
\label{sec:boundary-contacts-v157}
Let $R\subset M$ be a regular pencil and $\Lambda=\Pj(R)$. Its map
to $P$ lifts uniquely to a morphism
$\widetilde\iota_R:\Lambda\to\mathrm{CQ}(V)$: it lifts on the
nonsingular open, and properness extends it uniquely across the missing
points of the smooth curve. Put
$D_r(R)=\widetilde\iota_R^*E_r$. These are effective divisors because
the generic member is nonsingular.

\begin{theorem}[Contact-divisor reconstruction of spectral data]
\label{thm:contact-Segre-v157}
For every $h\geq1$, every regular pencil, and $0\leq p\leq nh$,
\begin{equation}\label{eq:contact-ideals-v157}
 J_p|_\Lambda=\OO_\Lambda\left(-\sum_{r=1}^{n-1}(p-hr)_+D_r(R)\right).
\end{equation}
At a point $x\in\Lambda$, let $0=e_1\leq\cdots\leq e_n$ be the
symmetric Smith exponents of the local pencil matrix. Then
\begin{equation}\label{eq:contact-exponents-v157}
 \operatorname{mult}_xD_r(R)=e_{r+1}-e_r\quad(1\leq r<n).
\end{equation}
Consequently the supported contact divisors determine every elementary
divisor, hence the full Segre symbol. If
$v_p(x)=\length(\OO_{\Lambda,x}/J_p)$ and $v_0(x)=0$, then
\begin{equation}\label{eq:contact-second-difference-v157}
 \operatorname{mult}_xD_r(R)=
 \frac{v_{h(r+1)}(x)-2v_{hr}(x)+v_{h(r-1)}(x)}{h}.
\end{equation}
\end{theorem}
\begin{proof}
Pull back \eqref{eq:power-boundary-v157}; equality of extended ideals
under composition gives \eqref{eq:contact-ideals-v157}.
The pencil matrix is nonzero at every point, since $R$ is a genuine
two-dimensional subspace; hence its entries generate the unit ideal
locally and $e_1=0$. Over $\C[[t]]$ symmetric elimination gives a
congruence to $\diag(t^{e_1},\ldots,t^{e_n})$, absorbing unit factors
by square roots. Its unique lift has, in
\eqref{eq:CQ-chart-v157}, $\tau_r=t^{e_{r+1}-e_r}$ up to a unit.
This proves \eqref{eq:contact-exponents-v157}; faithful flatness of
completion gives the local divisor statement. Taking valuations in
\eqref{eq:contact-ideals-v157} and second differences at multiples of
$h$ proves \eqref{eq:contact-second-difference-v157}.
Recover $e_i$ by summing the first $i-1$ contact multiplicities. The
positive exponents, with their points on $\Lambda$, are precisely
the elementary-divisor partitions of the spectral cokernel.
\end{proof}

\begin{example}[A boundary point and its contact order are distinct data]
\label{ex:contact-profile-v157}
For a symmetric germ with exponents $(0,1,3)$ and $h=2$, the valuations
of $J_1,\ldots,J_6$ are $(0,0,1,2,5,8)$.
Its two boundary contacts are $1$ and $2$. The germs
$\diag(1,t,t^3)$ and $\diag(1,t^2,t^5)$ reach the same full-flag
complete-quadric point but have contacts $(1,2)$ and $(2,3)$.
The compactification retains the flag; the arc with its divisorial
contacts retains the elementary exponents. These assertions must not
be conflated. For singular pencils the universal ideal theorem still
holds and detects all Fitting schemes and the normal rank; Fitting
ideals alone do not determine Kronecker minimal indices.
\end{example}

\section{A projective incidence compactification for pencils}
\label{sec:pencil-compactification-v157}
The complete-quadric graph concerns members of pencils. There is also
a projective parameter space carrying both the original finite failure
algebra and flat limits of their resolved member curves. This construction
keeps the pencil and does not identify distinct pencils at a common
full-coordinate-change limit.

\begin{proposition}[Resolved pencil curves and the failure family]
\label{prop:pencil-compactification-v157}
Let $n\geq3$, $G=\Gr(2,M)$, and let $G^\circ$ be the open set of
regular pencils whose lines avoid the rank-at-most-$(n-2)$ locus in
$P$. There is a normal projective $\operatorname{GL}(V)$-equivariant
scheme $\widehat G$ with a projective birational morphism
$\rho:\widehat G\to G$, an isomorphism over $G^\circ$, and a flat
family of subschemes
$\mathcal C\subseteq\widehat G\times\mathrm{CQ}(V)$ with the following
properties. Over $G^\circ$ it is exactly the lifted pencil line. Every
fibre over $R$ lies scheme-theoretically in $\pi^{-1}(\Pj(R))$.
All the resolved power maps restrict to $\mathcal C$, for every $h$.
The sharp finite failure algebra on $G$ pulls back to a finite locally
free algebra on $\widehat G$, with its original multiplication.
\end{proposition}
\begin{proof}
The rank-at-most-$(n-2)$ locus has codimension three in $P$, so a general
line avoids it; a general pencil is regular. Thus $G^\circ$ is nonempty
and dense. Over that open, $\pi$ is an isomorphism along every pencil
line and these lifts form a flat family of embedded projective lines.

For $1\leq q<n$, let $H_q$ be the pullback to $\mathrm{CQ}(V)$ of the
hyperplane bundle for its $q$th exterior-power projection, taking $q=1$
to mean $P$. The bundle $L=\bigotimes_{q=1}^{n-1}H_q$ is very ample,
by the graph embedding and Segre embedding. On a lifted line over
$G^\circ$, $\deg H_q=q$, since its $q$-minors have no common zero.
The Hilbert polynomial with respect to $L$ is therefore
$\binom n2 m+1$. The lifted family defines a morphism
\[
 G^\circ\longrightarrow
       \Hilb^{\binom n2 m+1}(\mathrm{CQ}(V)).
\]
Take the reduced closure of its graph in $G$ times this projective
Hilbert scheme, and then its normalization; this is $\widehat G$.
Normalization is finite here, so the result is projective and normal.
The graph is already closed over $G^\circ$ and that open is smooth,
so $\rho$ is an isomorphism there. Equivariance of the family and the
uniqueness of normalization give the stated group action.

Pull back the universal Hilbert family to obtain $\mathcal C$; it is
flat. To verify the asserted scheme-theoretic incidence, pull back
the linear equations for the universal pencil line to $\mathcal C$.
They vanish over the dense open $G^\circ$. On an affine open of the
integral base $\widehat G$, flatness makes the coordinate modules of
$\mathcal C$ torsion free over the base. A section which becomes zero
at the generic point is therefore zero. Applied locally to these
linear equations, this proves the incidence on the whole family.
The power maps are everywhere defined on $\mathrm{CQ}(V)$ by
Theorem~\ref{thm:power-graph-v157}, so their restrictions exist.
Finally pull back the finite locally free algebra of
Corollary~\ref{cor:finite-flat-family} along $\rho$; arbitrary base
change preserves its rank and multiplication.
\end{proof}

This is an incidence compactification over the Grassmannian, with its
native congruence action, not a properness assertion for the quotient
stack by that action. A limiting Hilbert fibre may have extra or
nonreduced components; no classification of all such fibres is used.
The construction is independent of $h$: its target, polarization and
Hilbert polynomial were specified using the $h$-independent complete
quadric graph. The unmarked inverse in the companion reconstruction
paper recovers the underlying pencil, so the incidence construction is
compatible with that inverse rather than a replacement of it.

\section{Singularities of the universal power-zero schemes}
\label{sec:power-multiplier-v157}
There is an explicit further consequence for all the universal ideals,
which we record separately from the geometry of the full reciprocal
normalization fibre. Multiplier ideals are taken on the affine smooth
space $M$, and $I_{r+1}^{(a)}$ denotes symbolic power, with exponent
zero interpreted as the unit ideal.

\begin{corollary}[Multiplier ideals and threshold]
\label{cor:power-multiplier-v157}
For $1\leq p\leq nh$, put $w_r=(p-hr)_+$ and
$c_r=\binom{n-r+1}{2}$ for $0\leq r<n$. For every real $b\geq0$,
\begin{align}
 \mathcal J(M,b\cdot J_p)
   &=\bigcap_{r=0}^{n-1}
       I_{r+1}^{(\max\{0,\lfloor b w_r\rfloor+1-c_r\})},
                         \label{eq:power-multiplier-v157}\\
 \operatorname{lct}_M(J_p)
   &=\min_{\substack{0\leq r<n\\p>hr}}
                    \frac{\binom{n-r+1}{2}}{p-hr}.
                         \label{eq:power-lct-v157}
\end{align}
\end{corollary}
\begin{proof}
Blow up the affine origin and then the strict rank loci. This is the
affine complete-quadric resolution, whose charts are
\eqref{eq:CQ-chart-v157} with one additional scalar coordinate.
Let $E_0$ be the divisor over the origin. The same minor calculation
now includes $r=0$, and the total transform of $J_p$ is
$\sum_{r=0}^{n-1}w_rE_r$. The discrepancy along $E_r$ is $c_r-1$:
over the generic rank-$r$ point the earlier blow-ups are isomorphisms,
and the smooth rank-$r$ locus has codimension $c_r$; for $r=n-1$
this is the strict determinant divisor and the discrepancy is zero.
All boundary divisors have normal crossings, so this is a log resolution.

The defining pushforward formula for the multiplier ideal requires
$\operatorname{ord}_{E_r}(f)\geq\lfloor b w_r\rfloor+1-c_r$.
At the generic rank-$r$ point this valuation is the order in the
regular local ring with maximal ideal $I_{r+1}$, so the condition on
polynomials is exactly membership in the indicated symbolic power.
Intersecting the conditions proves the first formula; the first value
at which one becomes nontrivial proves the second. Equivalently this
is the specialization, via Theorem~\ref{thm:universal-power-v157}, of
the classical symmetric-determinantal multiplier-ideal calculation
\cite[\S2.4 and Theorem~4.8]{HenriquesVarbaro2016}.
\end{proof}

The ideal formula, graph identification and contact reconstruction are
statements about the entire universal coefficient space. They do not
supply the conductor of every multivariate common-divisor image or the
Hilbert function of every full $F_{g,k}(W)$. Those objects remain in
the preserved incidence theory; the present results connect their exact
quadratic sections to a native global boundary and its spectral schemes.
